\providecommand{\Res}[2]{\qty[round-mode=figures,round-precision=3]{\fpeval{#1}}{#2}}
\providecommand{\Val}[1]{\num[round-mode=figures,round-precision=4]{\fpeval{#1}}}
\providecommand{\ResSci}[2]{\qty[round-mode=figures,round-precision=3,exponent-mode=scientific]{\fpeval{#1}}{#2}}
\subsubsection*{Question a) Position qualitative}
À l'extérieur du segment, les deux forces ont le même sens (les deux repoussent ou les deux attirent $Q$) : elles ne peuvent pas s'annuler. Entre les charges, la force de $q_1$ et celle de $q_2$ sont de sens opposés ; elles peuvent s'annuler. Le point d'équilibre est plus près de la plus petite charge.

\subsubsection*{Question b) Position d'équilibre}
\begin{enumerate}[label=\arabic*)]
    \item \textbf{Variables connues et recherchées :} $q_1 = \qty{\v{q1}e-6}{C}$, $q_2 = \qty{\v{q2}e-6}{C}$, $d = \qty{\vsi{d}}{m}$ ; variable recherchée : $x$ (distance à partir de $q_1$)
    \item \textbf{Équations pertinentes :}
    \[ \frac{k q_1 Q}{x^2} = \frac{k q_2 Q}{(d - x)^2} \]
    \item \textbf{Isolation de la variable recherchée :} en prenant la racine carrée (avec $0 < x < d$) :
    \[ \frac{d - x}{x} = \sqrt{\frac{q_2}{q_1}} \quad \Rightarrow \quad x = \frac{d}{1 + \sqrt{q_2 / q_1}} \]
    \item \textbf{Résolution numérique :}
    \[ x = \frac{\num{\vsi{d}}}{1 + \sqrt{\dfrac{\num{\v{q2}}}{\num{\v{q1}}}}} = \Res{\vsi{d}/(1 + sqrt(\v{q2}/\v{q1}))}{m} \]
    \item \textbf{Réponse :} il faut placer $Q$ à $x = \mathbf{\Res{100*\vsi{d}/(1 + sqrt(\v{q2}/\v{q1}))}{cm}}$ de $q_1$. Ce résultat ne dépend ni du signe ni de la valeur de $Q$.
\end{enumerate}

\subsubsection*{Question c) Force au milieu}
\begin{enumerate}[label=\arabic*)]
    \item \textbf{Variables connues et recherchées :} $Q = \qty{\v{Q}e-9}{C}$, distance de chaque charge $d/2 = \qty{\fpeval{\vsi{d}/2}}{m}$ ; variable recherchée : $F_{\text{rés}}$
    \item \textbf{Équations pertinentes :} les deux forces sont opposées :
    \[ F_{\text{rés}} = \frac{kQ\,|q_1 - q_2|}{(d/2)^2} \]
    \item \textbf{Isolation de la variable recherchée :} l'équation est déjà isolée.
    \item \textbf{Résolution numérique :}
    \[ F_{\text{rés}} = \frac{\num{8.99e9} \cdot \num{\v{Q}e-9} \cdot \num{\fpeval{abs(\v{q1}-\v{q2})}e-6}}{(\num{\fpeval{\vsi{d}/2}})^2} = \Res{8.99e9*\v{Q}*1e-9*abs(\v{q1}-\v{q2})*1e-6/(\vsi{d}/2)^2}{N} \]
    \item \textbf{Réponse :} $F_{\text{rés}} = \mathbf{\Res{8.99e9*\v{Q}*1e-9*abs(\v{q1}-\v{q2})*1e-6/(\vsi{d}/2)^2}{N}}$, dirigée vers la plus petite des deux charges.
\end{enumerate}
